\documentclass[a4paper]{article}
\usepackage[utf8]{inputenc}
\usepackage[T1]{fontenc}
\usepackage[french]{babel}
\usepackage{graphicx}
\usepackage{hyperref}
\usepackage{amsmath,amssymb}
\usepackage{geometry}
\geometry{margin=2.5cm}
\title{Projet d'électronique E3 EISC - Partie 1}
\author{Votre Nom}
\date{2025-2026}
\begin{document}
\maketitle
\begin{center}
  \includegraphics[width=0.5\textwidth]{Pictures/100000010000019F000000596F16D858.png} \\
  \textbf{Année Universitaire 2025-2026}
\end{center}

\section*{Introduction à LTspice}
LTspice est un simulateur circuit basé sur le standard \textbf{SPICE}, développé par la société Analog Devices. Cet outil logiciel gratuit n'a pas de limitation sur le nombre de nœuds ou de composants, c'est un standard de l'industrie pour la simulation de circuits analogiques. Pour analyser ou concevoir un circuit, c'est un outil indispensable qui permet de passer d'un schéma théorique à une visualisation précise des signaux électriques (tensions, courants, spectres de fréquence) avant d'envisager un premier prototype.

Pour télécharger le logiciel LTspice, rendez-vous sur :
\url{https://www.analog.com/en/resources/design-tools-and-calculators/ltspice-simulator.html}

Pour la prise en main du logiciel, rendez-vous sur le LT wiki :
\url{https://ltwiki.org/}

\section*{Introduction du projet}
L'objectif du projet est de réaliser à l'aide du logiciel LTspice les études statique et dynamique de différents montages intégrant un transistor bipolaire de type NPN 2SC5876U3 de la société Rohm semiconductor. Toutes les informations sur le composant sont disponibles au lien suivant :
\url{https://www.rohm.com/products/transistors/bipolar-transistors/automotive-bipolar-transistors/2sc5876u3hzg-product#productDetail}

\begin{center}
  \begin{tabular}{c c}
    \includegraphics[width=0.2\textwidth]{Pictures/100000000000005A000000640960B587.png} &
    \includegraphics[width=0.2\textwidth]{Pictures/10000001000000BD0000008291F1FE3D.png} \\
    Symbole du transistor NPN & Transistor 2SC5876U3 : élément LTspice \\
  \end{tabular}
\end{center}


\section*{PARTIE 1 : Caractérisation du transistor}
On souhaite établir les caractéristiques du transistor 2SC5876U3 en émetteur commun présentées sur la figure suivante :

\begin{figure}[h!]
  \centering
  \includegraphics[width=0.6\textwidth]{Pictures/10000001000001C2000000F975CC44B1.png}
  \caption{Montage d'analyse des caractéristiques du transistor}
\end{figure}

On se propose donc d'étudier chacune des caractéristiques du transistor :
\begin{itemize}
  \item $I_C = f(I_B)$
  \item $I_C = f(V_{CE})$
  \item $V_{BE} = f(I_{B})$
  \item $V_{BE} = f(V_{CE})$
\end{itemize}

\textbf{Note :} Pour $I_C = f(V_{CE})$ on choisira 3 à 4 valeurs de $I_B$ dans la gamme de compatibilité du transistor.

Pour cela, une possibilité est l'utilisation du montage de la figure 1 et du mode de simulation \texttt{.DC sweep}.

\section*{Q1. Démarche de simulation et analyse des caractéristiques}
En détaillant la démarche de simulation, il s'agit de tracer chacune des caractéristiques du transistor et de repérer les zones de fonctionnement. À partir de l'analyse des courbes obtenues, il faut vérifier l'adéquation entre les résultats de simulation, la documentation technique et les paramètres du modèle SPICE du composant.

Pour tracer les courbes suivantes $I_C=f(I_B)$, $I_C=f(V_{CE})$, $V_{BE}=f(I_B)$, $V_{BE}=f(V_{CE})$, nous commençons par créer un projet sur LTspice et réaliser le montage de la Figure 1.

\subsection*{Caractéristique Gain Courant : $I_C=f(I_B)$}
Pour cette caractéristique, il faut :
\begin{itemize}
  \item Effectuer un DC sweep sur la source de tension $V_1$ qui alimente $V_{BB}$ (contrôle $I_B$ via $R_B$)
  \item Utiliser une source continue simple pour $V_2$
\end{itemize}
\begin{figure}[h!]
  \centering
  \includegraphics[width=0.5\textwidth]{Pictures/100000010000021C000001E40C231913.png}
  \caption{Montage LTspice pour $I_C=f(I_B)$}
\end{figure}

On obtient le montage suivant :
\begin{figure}[h!]
  \centering
  \includegraphics[width=0.6\textwidth]{Pictures/100000010000036200000224DC22A322.png}
  \caption{Montage complet pour la caractéristique $I_C=f(I_B)$}
\end{figure}

Pour tracer la caractéristique $I_C=f(I_B)$, il faut tracer $I_C(Q1)$ ou $I(R_C)$ sur l’axe Y et $I(R_B)$ sur l’axe X. Ce qui donne :
\begin{figure}[h!]
  \centering
  \includegraphics[width=0.9\textwidth]{Pictures/100000000000077C0000037DB4C056F4.png}
  \caption{Exemple de courbe obtenue pour $I_C=f(I_B)$}
\end{figure}

\subsection*{Caractéristique $I_C=f(V_{CE})$}
Pour cela, il faut faire varier la tension entre collecteur et émetteur, ce qu’on peut faire en faisant varier $V_{CC}$ alimenté par $V_2$.

\textbf{Note :} Pour $I_C=f(V_{CE})$ on choisira 3 à 4 valeurs de $I_B$ dans la gamme de compatibilité du transistor.

Pour choisir la plage de courant de base ($I_B$), il faut s’assurer que le point de fonctionnement du transistor ne dépasse jamais les limites destructives définies par le constructeur.

\section*{Q2. Gain en courant $\beta = f(I_C)$}
Tracer l’évolution du gain en courant $\beta = f(I_C)$ (pour $V_{CE}$ fixé) et comparer avec la documentation technique.

\section*{Q3. Gain en tension}
À partir de la caractéristique de transfert en tension $V_{CE}= f(V_{BE})$, tracer l’évolution du gain en tension en fonction du point de polarisation.

\section*{PARTIE 2 : Analyse de cas / conception d’un amplificateur}
Vous décidez de vous lancer dans la conception de l’amplificateur à base du transistor 2SC5876U3 caractérisé dans la partie précédente. Vous travaillez dans un premier temps à vide (sans charge), c'est-à-dire sans connecter le haut-parleur.

Vous testez les différents montages de la figure 2. La simulation devrait aider à comprendre le fonctionnement de ces montages et permettre de comparer avec la théorie.

% Figure 2 : Etude de l’amplificateur à vide (à insérer ici)

\section*{Q5. Simulation des montages de la figure 2}
Pour chaque montage de la figure 2, effectuer la simulation du montage et indiquer si la structure permet d’amplifier correctement le signal issu du microphone. Si l’amplification est correcte, on caractérisera l’amplificateur en termes d’impédances (entrée et sortie) et de gains sur l’ensemble de sa bande passante.

Vous étudiez ensuite l’amplificateur en charge en connectant le haut-parleur. Vous simulez les différents montages indiqués sur la figure 3.

% Figure 3 : Amplificateur en charge (à insérer ici)

\section*{Q6. Simulation des montages de la figure 3}
Pour chaque montage de la figure 3, indiquer si la structure permet d’amplifier correctement le signal issu du microphone. Si l’amplification est correcte, on caractérisera l’amplificateur en termes d’impédances (entrée et sortie) et de gains sur l’ensemble de sa bande passante.

Les résultats précédents n’étant pas satisfaisants, vous décidez de réaliser la structure à deux étages présentée sur la figure 4. Les deux transistors utilisés sont identiques.

% Figure 4 : Amplificateur audio à deux étages en charge (à insérer ici)

\section*{Q7. Simulation du montage de la figure 4}
Simuler le montage de la figure 4 et caractériser l’amplificateur en termes d’impédances (entrée et sortie) et de gains sur l’ensemble de sa bande passante. Donner vos conclusions sur le rôle du second étage.

% --- Fin du document ---
\end{document}
